\section{Conclusão}\label{sec:conclusion}

Este trabalho apresentou um fluxo demonstrativo para validação automatizada de documentos acadêmicos produzidos em \LaTeX{}. A proposta separou preparação, compilação, verificações automáticas e revisão visual, mostrando como propriedades objetivas do PDF podem ser verificadas sem atribuir à automação decisões que continuam dependentes de análise humana.

O exemplo também demonstrou uma forma de organizar o projeto acadêmico: o arquivo principal concentra metadados e ordem dos elementos, enquanto capítulos, resumo, abstract, referências e materiais complementares permanecem em arquivos específicos. Figuras, tabelas, equações e código aparecem no contexto metodológico em que são utilizados.

Como limitação, os resultados apresentados são sintéticos e servem apenas para mostrar a estrutura do documento. Em um TCC real, objetivos, método, dados, resultados e conclusões devem decorrer da pesquisa efetivamente realizada. A documentação completa dos comandos da classe permanece separada deste texto para que o PDF do aluno continue sendo um trabalho acadêmico, e não um manual do template.
